Independent tabular Q-learning pricing agents in a repeated Bertrand duopoly are known to converge on supracompetitive prices without explicit communication \citep{calvano2020}. The literature has so far studied this phenomenon under idealised informational and timing assumptions: instantaneous and noiseless observation of the rival's price, and simultaneous price-setting each period. This paper relaxes those assumptions one at a time and then in combination. Three frictions are introduced into the otherwise unchanged \citet{calvano2020} environment: (i) \emph{observation latency}, where each agent observes the rival's price with a delay of $L$ periods; (ii) \emph{noisy monitoring}, where the rival's observed price is corrupted by Gaussian measurement error of standard deviation $\sigma$; and (iii) \emph{asynchronous actions} in the Maskin--Tirole sense, where only one agent revises its price per $T$-period block. The single-friction sweeps reveal three qualitatively different shapes of the collusion-vs-friction relationship: latency exhibits a sharp cliff at $L=1$ followed by a residual collusion floor of about $\Delta\approx 0.16$; noise produces a smooth monotone decline that reaches Nash by $\sigma=5$; asynchrony reduces $\Delta$ from $0.84$ to $\approx 0.40$ at $T=1$ and then plateaus. A $2\times 2\times 2$ factorial experiment, replicated at $250$, $500$ and $1000$ sessions, examines whether the frictions interact. Two robust interaction patterns emerge: latency and noise compound super-multiplicatively, and asynchrony partially \emph{rescues} collusion when imposed on top of either of the other two frictions. The latter result is the opposite-sign of the main effect of asynchrony alone, and suggests that mandatory staggered-pricing rules and information-degradation rules are policy substitutes rather than complements.
